\documentclass[11pt]{article}
\usepackage[ngerman]{babel}

\usepackage[T1]{fontenc}
\usepackage[utf8]{inputenc}
\usepackage{lmodern}
\usepackage{amsmath,amssymb,mathtools,array,booktabs,pdflscape,adjustbox,enumitem}
\usepackage{geometry}
\usepackage{microtype}
\geometry{a4paper,margin=1.45cm}
\setlength{\parindent}{1.25em}
\setlength{\parskip}{0pt}
\makeatletter
\renewcommand\footnotesize{\@setfontsize\footnotesize{7.5}{8.5}\selectfont}
\makeatother
\allowdisplaybreaks
\setlength{\emergencystretch}{16em}
\sloppy
\hbadness=10000
\vbadness=10000
\hfuzz=3pt
\vfuzz=10pt
\raggedbottom
\providecommand{\frakS}{\mathfrak S}
\providecommand{\frakM}{\mathfrak M}
\providecommand{\frakG}{\mathfrak G}
\providecommand{\frakm}{\mathfrak m}
\providecommand{\frakn}{\mathfrak n}
\providecommand{\frakp}{\mathfrak p}
\providecommand{\frakq}{\mathfrak q}
\providecommand{\frakr}{\mathfrak r}
\providecommand{\fraka}{\mathfrak a}

\providecommand{\frakb}{\mathfrak b}
\providecommand{\frakc}{\mathfrak c}
\providecommand{\barfraka}{\overline{\mathfrak a}}
\providecommand{\barfrakb}{\overline{\mathfrak b}}
\providecommand{\barfrakc}{\overline{\mathfrak c}}
\providecommand{\frakb}{\mathfrak b}
\providecommand{\frakc}{\mathfrak c}
\providecommand{\frakd}{\mathfrak d}
\providecommand{\frake}{\mathfrak e}
\providecommand{\frakg}{\mathfrak g}
\providecommand{\fraks}{\mathfrak s}
\providecommand{\frako}{\mathfrak o}
\providecommand{\frf}{\mathfrak f}
\providecommand{\frK}{\mathfrak K}
\providecommand{\frR}{\mathfrak R}
\providecommand{\calP}{\mathcal P}
\providecommand{\calR}{\mathcal R}
\providecommand{\calN}{\mathcal N}
\providecommand{\calG}{\mathcal G}
\providecommand{\calH}{\mathcal H}
\providecommand{\barfrakm}{\overline{\mathfrak m}}
\providecommand{\barr}{\overline}
\providecommand{\dd}{\mathrm d}
\providecommand{\tagg}[1]{\tag{#1}}
\begin{document}
\setcounter{footnote}{0}

\section*{24. Eliminationstheorie und allgemeine Idealtheorie}
\begin{center}
Math. Ann. 90 (1923), S. 229--261
\end{center}

Es handelt sich im folgenden um eine Einordnung der Eliminationstheorie -- in der arithmetischen Fassung, die ich der Hentzeltschen Darstellung gegeben habe\footnote{Hentzelt, Kurt: Zur Theorie der Polynomideale und Resultanten. Bearbeitet von E. Noether. Math. Ann. 88 (1922), S. 53--79; zitiert H.-N.} und die sich als eine Idealtheorie im Polynombereich bezeichnen läßt -- in die allgemeine Idealtheorie,\footnote{Noether, E.: Idealtheorie in Ringbereichen, Math. Ann. 83 (1921), S. 24--66; zitiert Idealtheorie.} zugleich um eine Neubegründung der auf die Nullstellen bezüglichen Teile. Die Grundbegriffe beider Theorien sind in § 1 kurz zusammengestellt, ebenso die der Körpertheorie,\footnote{Steinitz, E.: Algebraische Theorie der Körper, J. f. M. 137 (1910), S. 167--309; zitiert Steinitz.} so daß ich mich hier darauf beziehen kann.

Einordnung und Neubegründung gruppieren sich um vier Fragestellungen: Arithmetische Fassung des Dimensionsbegriffes; Theorie der Nullstellen; Zusammenhang der Zerlegungssätze für Normen und Elementarteiler mit denen der allgemeinen Idealtheorie; Charakterisierung der Primideale und Primärideale durch Norm und Elementarteilerform.

Die arithmetische Fassung des \emph{Dimensionsbegriffes} (§ 4) geschieht durch Ketten von Primidealen, das arithmetische Äquivalent dafür, daß es auf Flächen Kurven, auf Kurven Punkte gibt. Diese arithmetische Fassung, die mit der Parameterdefinition der Eliminationstheorie als identisch nachgewiesen wird, läßt sich mit leichter Modifikation auf beliebige Ringbereiche übertragen und gibt dort zusammen mit der Übertragung gewisser Teile der übrigen Fragestellungen eine genaue Einsicht in den Aufbau der Ideale, wie ich an anderer Stelle zeigen werde.

Die Frage der \emph{Nullstellen} wird bei H.-N. wie üblich für das gegebene Ideal direkt in Angriff genommen, und zwar durch Zurückgehen auf sukzessive Elimination, wobei der Charakter des Verifizierens nicht ganz vermieden wird. Im Gegensatz dazu steht der für Polynome einer Veränderlichen stets begangene Weg; man stellt das gegebene Polynom als Produkt von Potenzen von Primfunktionen (Primärfunktionen) dar und konstruiert für jede einzelne Primfunktion den Nullstellenkörper als einen dem Restklassenkörper isomorphen Körper. Der Grad der Primfunktion gibt den Grad dieses Körpers; der Grad der Primärfunktion aber, deren Nullstellen mit denen der Primfunktion übereinstimmen, gibt den Grad des Ringes der Restklassen nach dieser Primärfunktion. Geht man bei jeder Primfunktion zum Galoisschen Körper über, so kommt die Zerlegung in Linearfaktoren; und damit ist für das ursprünglich gegebene Polynom die Frage der Nullstellen, Zerlegung in Linearfaktoren und Multiplizität erledigt. Ganz entsprechend wird hier (§ 3) das System der Restklassen nach einem Primideal durch Quotientenbildung zum Restklassenkörper erweitert und ein dazu isomorpher Nullstellenkörper konstruiert. Dieser Nullstellenkörper enthält einen durch Adjunktion von Unbestimmten -- etwa $x_{i+1},\ldots,x_n$ -- erzeugten Teilkörper; der Grad der Norm gibt den Grad in bezug auf diesen Teilkörper; zugleich ergibt sich durch Übergang zum Galoisschen Körper die Zerlegung der Elementarteilerform und damit der Norm in Linearfaktoren. Für das zugehörige Primärideal aber sind die Nullstellen die gleichen; die Zerlegung ist ebenfalls mitgegeben, da die Norm eine Potenz der Elementarteilerform des Primideals wird (§ 2); der Grad der Norm gibt den Grad des Ringes der Restklassen nach dem Primärideal in bezug auf den aus den Unbestimmten abgeleiteten Teilkörper.

Damit ist aber die Frage der Nullstellen eines beliebigen Ideals zurückgeführt auf den Zusammenhang der Zerlegungssätze für Normen und Elementarteiler mit denen der allgemeinen Idealtheorie (§ 5). Der Tatsache, daß die Nullstellen des Ideals sich aus denen seiner einzelnen zugehörigen Primideale zusammensetzen, entspricht die Zerlegung der Norm in größte primäre Faktoren, die gleich Potenzen der Elementarteilerform der einzelnen zugehörigen Primideale werden;\footnote{Dieses Parallelgehen der Eliminationstheorie mit der allgemeinen Idealtheorie ist bei der Kroneckerschen Eliminationstheorie nicht erfüllt; vgl. dazu H.-N., Anm. *).} damit ist die Zerlegung in Linearfaktoren für die Norm allgemein geleistet. Die Multiplizität aber findet bei H.-N. ihre Deutung vermöge der Grundideale und ihrer Komponenten; es wird der Grad eines größten primären Faktors der Norm gleich der Anzahl linear unabhängiger Restklassen -- bei Adjunktion von $x_{i+1},\ldots,x_n$ -- des Grundideals $(i-1)$-ter Stufe nach einer Komponente des Grundideals $i$-ter Stufe. Nun wird das Grundideal $(i-1)$-ter Stufe als identisch nachgewiesen mit der isolierten Komponente der Zerlegung, die eindeutig definiert ist durch alle Primideale von höherer als der $(n-i)$-ten Dimension; eine Komponente des Grundideals $i$-ter Stufe aber als identisch mit der isolierten Komponente der Zerlegung, die definiert ist durch das dem Faktor der Norm entsprechende Primideal und alle Primideale höherer Dimension. Und der Grad des entsprechenden Faktors der Norm wird -- bei Adjunktion von $x_{i+1},\ldots,x_n$ -- gleich dem Grad desjenigen Teilringes der Restklassen dieser Komponente, der nur aus Klassen des Grundideals $(i-1)$-ter Stufe besteht.

Die Charakterisierung der Primideale und eigentlichen Primärideale (§ 6) ergibt sich als direkte Folgerung der Betrachtung der Restklassenkörper. Ist der ursprüngliche Koeffizientenbereich ein vollkommener Körper, so entsprechen den Primidealen als Norm und Elementarteilerform Primfunktionen, den eigentlichen Primäridealen eigentliche Primärfunktionen, und umgekehrt. Bei unvollkommenem Körper als Koeffizientenbereich lassen sich nur Bedingungen angeben, die entweder notwendig oder hinreichend sind; es kann, wie an Beispielen gezeigt wird, die Norm eines Primideals eigentlich primär werden, die Elementarteilerform eines eigentlichen Primärideals eine Primfunktion. Genauer wird hier bei Charakteristik $p$ die Norm eines Primideals die $p^f$-te Potenz ($f>0$) seiner Elementarteilerform; während bei eigentlichen Primäridealen, deren Elementarteilerform eine Primfunktion ist, die Norm eine beliebige Potenz werden kann, wie wieder Beispiele zeigen. Es werden schließlich noch absolute Primideale betrachtet (§ 7), d. h. solche, die Primideale bleiben bei Erweiterung des Koeffizientenbereiches zu einem algebraisch abgeschlossenen Körper. Bei absoluten Primidealen mit Koeffizienten aus einem (endlichen) algebraischen Zahlkörper führt die Charakterisierung zur Übertragung eines zuerst von A. Ostrowski für absolute Primfunktionen aufgestellten Satzes: Es bleibt die Eigenschaft, Primideal zu sein, erhalten modulo jedes Primideals des Zahlkörpers, höchstens mit endlich vielen Ausnahmen.

Es sei noch auf die Analogie der letzten Fragestellung mit der Idealtheorie in algebraischen Zahlkörpern hingewiesen. Als Elementarteilerform eines solchen Ideals ist die kleinste durch das Ideal teilbare ganze rationale Zahl aufzufassen (Anm. 10), die also für ein Primideal stets Primzahl wird, während umgekehrt ein Ideal Primideal wird, sobald seine Norm Primzahl ist; auch die Beweise laufen ganz parallel. Die oben angeführten Beispiele zeigen also, daß bei unvollkommenem Körper das Analogon zu Primidealen höheren Grades und zu Verzweigungsidealen auftritt, während es bei vollkommenem Körper nur Primideale ersten Grades und keine Verzweigungsideale gibt.

\subsection*{§ 1. Grundbegriffe der Eliminationstheorie, allgemeinen Idealtheorie und Körpertheorie}

\textbf{1. Der Bereich.} Es bezeichne $\frakS$ den Integritätsbereich aller Polynome in $y_1,\ldots,y_n$ mit Koeffizienten aus einem abstrakt definierten Körper $P$. Die $y$ seien einer Transformation mit Unbestimmten $u_{\mu\nu}$ als Koeffizienten unterworfen:\footnote{Die Transformation ist im Hinblick auf § 3 usw. etwas allgemeiner als bei H.-N., wodurch alle dortigen Resultate um so mehr erhalten bleiben.}
\begin{align}
 y_1&=u_{11}x_1+\cdots+u_{1n}x_n,\quad \ldots,\quad
 y_n=u_{n1}x_1+\cdots+u_{nn}x_n, \tag{1}
\end{align}
mit der Umkehrung
\begin{align}
 x_1&=t_{11}y_1+\cdots+t_{n1}y_n,\quad \ldots,\quad
 x_n=t_{1n}y_1+\cdots+t_{nn}y_n, \tag{1'}
\end{align}
und die $u_{\mu\nu}$, oder -- was wegen der gegenseitig rationalen Ausdrückbarkeit damit gleichbedeutend ist -- die $t_{\mu\nu}$ zu $P$ adjungiert, so daß der Koeffizientenbereich $P(u)=P(t)$ entsteht; $\frakS'$ bezeichne den Integritätsbereich aller Polynome in $x_1,\ldots,x_n$ mit Koeffizienten aus $P(u)$. Die Bereiche $\frakS$ und $\frakS'$ liegen der Eliminationstheorie zugrunde; es werden Ideale $\fraka,\frakb,\ldots$ in $\frakS$ und $\fraka',\frakb',\ldots$ in $\frakS'$ betrachtet. Ein Ideal in einem beliebigen Ring ist dabei wie üblich dadurch definiert, daß es neben $\alpha$ und $\beta$ auch $\alpha-\beta$ enthält, neben $\alpha$ auch $\lambda\alpha$, unter $\lambda$ ein beliebiges Ringelement verstanden; $\frako$ bedeute stets das aus allen Elementen bestehende Einheitsideal.

\textbf{2. Transformierte Ideale.} Ein Ideal $\frakm$ in $\frakS'$ heißt transformiert, wenn es aus einem Ideal $\barfrakm$ in $\frakS$ vermöge (1) bei Adjunktion der $u_{\mu\nu}$, oder $t_{\mu\nu}$, entsteht. $\frakm$ besteht also, wenn $\bar f(y)$ oder $\bar g(y)$ alle Polynome aus $\barfrakm$ durchlaufen, aus allen linearen Verbindungen
\[
 f(x)=\sum U_i(u)\bar f_i(y)=\sum U_i(u)f_i(x)
\]
oder auch
\[
 g(x)=\sum T_i(t)\bar g_i(y)=\sum T_i(t)g_i(x);
\]
hierin sind insbesondere alle $f_i(x)=\bar f_i(y)$ enthalten, oder, was auf dasselbe hinauskommt, alle $g_i(x)=\bar g_i(y)$. Da $f(x)$ bzw. $g(x)$ nur bis auf Größen aus $P(u)=P(t)$ festgelegt sind, dürfen die $U_i,T_i$ ohne Beschränkung der Allgemeinheit als Potenzprodukte der $u$ bzw. $t$ angenommen werden. $f(x),g(x)$ gehen wieder in Polynome aus $\frakm$ über, wenn $U,T$ durch irgendwelche anderen Potenzprodukte ersetzt werden, also insbesondere durch Vertauschen der Spalten von (1) bzw. der Zeilen von (1'); somit enthält $\frakm$ neben irgendeinem Polynom auch alle durch Vertauschung der $x$ daraus entstehenden, sofern nur in den von $u$ bzw. $t$ abhängenden Koeffizienten die entsprechenden Vertauschungen vorgenommen werden. Alle im folgenden auftretenden Ideale sollen -- wenn nicht ausdrücklich das Gegenteil bemerkt ist -- als transformiert angenommen werden. Nichttransformierte Ideale gehen durch Zusammensetzen von (1) mit
\[
 x_i=v_{i1}z_1+\cdots+v_{in}z_n
\]
in transformierte im Polynombereich der $z$ mit Koeffizienten aus $P(u,v)$ über, während dieses Zusammensetzen bei transformierten Idealen nur auf ein Ersetzen der $u_{\mu\nu}$ durch bilineare Verbindungen der $u,v$ herauskommt.

\textbf{3. Bezeichnung.} Unter $f^{(i)},g^{(i)},\ldots,a^{(i)},b^{(i)},\ldots$ seien durchweg Polynome in $\frakS'$ verstanden, die von $x_i,\ldots,x_{i-1}$ frei sind.

\textbf{4. Grundideale.} Jedem Ideal $\frakm$ sind $n$ Grundideale $\frakg_0,\frakg_1,\ldots,\frakg_{n-1}$ zugeordnet durch die Festsetzung: Das Grundideal $(i-1)$-ter Stufe $\frakg_{i-1}$ enthält alle und nur die Polynome $G(x)$, für die es ein -- im allgemeinen mit $G(x)$ sich änderndes -- Polynom $b^{(i)}\ne 0$ gibt, derart, daß $b^{(i)}G(x)\equiv 0\pmod{\frakm}$ wird. Aus der Existenz der Idealbasis folgt daraus auch die Existenz mindestens eines für alle $G(x)$ festen $B^{(i)}$, so daß also
\begin{equation}
 B^{(i)}\frakg_{i-1}\equiv0\pmod{\frakm},\qquad B^{(i)}\ne0 . \tag{2}
\end{equation}
Die Grundideale transformierter Ideale werden selbst transformierte Ideale (H.-N. Satz VI). Das Grundideal $i$-ter Stufe $\frakg_i$ ist auch definiert als das Ideal, in das $\frakm$ bei Adjunktion von $x_{i+1},\ldots,x_n$ zu $P(u)$ übergeht, wenn man sich -- was dann keine Beschränkung der Allgemeinheit bedeutet -- auf in $x_{i+1},\ldots,x_n$ ganze Polynome beschränkt; $\frakg_0$ wird stets gleich dem Einheitsideal $\frako$.

\textbf{5. Moduldarstellung der Ideale, Elementarteiler und Einzelnormen.} Faßt man jedes Polynom $f(x)$ auf als Linearform in den Potenzprodukten $\xi$ von $x_1,\ldots,x_{i-1}$, mit Polynomen $a^{(i)}$ als Koeffizienten, so geht das Ideal $\frakm$ über in einen Modul $\frakM_{i-1}$ aus Linearformen in bezug auf den Bereich der $a^{(i)}$; d. h. $\frakM_{i-1}$ enthält neben irgend zwei Linearformen auch ihre Differenz, neben einer Linearform $l(\xi)$ auch $a^{(i)}l(\xi)$ bei beliebigem $a^{(i)}$. Als Grundmodul eines solchen Moduls $\frakM$ wird die Gesamtheit der Linearformen $g(\xi)$ bezeichnet, für die $b^{(i)}g(\xi)\equiv0\pmod{\frakM}$, $b^{(i)}\ne0$ wird; das Grundideal $\frakg_{i-1}$ geht also bei der Moduldarstellung in den Grundmodul $\frakG_{i-1}$ von $\frakM_{i-1}$ über. $\frakM_{i-1}$ und $\frakG_{i-1}$ hängen von unendlich vielen Unbestimmten $\xi$ ab, derart, daß in jeder einzelnen Linearform nur endlich viele dieser Unbestimmten auftreten. $\frakM_{i-1}$ hat die charakteristische Eigenschaft, bei Adjunktion von $x_{i+1},\ldots,x_n$ zu $P(u)$ nur endlich viele von der Einheit verschiedene Elementarteiler zu besitzen; d. h. $\frakG_{i-1}$ und $\frakM_{i-1}$ lassen bei dieser Adjunktion eine Basisdarstellung zu -- unter $\zeta$ neue Unbestimmte verstanden, die mit den $\xi$ durch umkehrbar lineare, in $x_i$ ganze Transformation zusammenhängen, derart, daß in $\zeta_i=r_i(\xi)$ nur jeweils endlich viele dieser Unbestimmten auftreten --
\begin{align}
\frakG_{i-1}&=(\zeta_0,\zeta_1,\ldots,\zeta_\sigma,\zeta_{\sigma+1},\ldots,\zeta_{\sigma+\nu},\ldots), \notag\\
\frakM_{i-1}&=(E_0^{(i)}\zeta_0,E_1^{(i)}\zeta_1,\ldots,E_\sigma^{(i)}\zeta_\sigma,
\zeta_{\sigma+1},\ldots,\zeta_{\sigma+\nu},\ldots), \tag{3}
\end{align}
wo jedes $E_\mu^{(i)}$ im vorangehenden aufgeht (H.-N. (24) und (28)).

Der höchste Elementarteiler $E^{(i)}$ ist auch definiert als größter gemeinsamer Teiler -- im Polynomsinn -- aller in (2) auftretenden $B^{(i)}$; er kann als ganz und primitiv in $x_{i+1},\ldots,x_n$ angenommen werden und wird dann regulär in $x_i$, d. h. $E^{(i)}$ enthält den Term $x_i^\lambda$ mit von Null verschiedenem Koeffizienten, wenn es vom Grade $\lambda$ in den $x$ ist.

Das Produkt $R^{(i)}$ aller in (3) auftretenden Elementarteiler wird als Norm von $\frakG_{i-1}$ nach $\frakM_{i-1}$ oder auch als Einzelnorm des Ideals $\frakm$ bezeichnet; in Zeichen und unter Berücksichtigung, daß $\frakm$ bei der Adjunktion von $x_{i+1},\ldots,x_n$ in $\frakg_i$ übergeht:
\begin{equation}
 R^{(i)}=E_0^{(i)}E_1^{(i)}\cdots E_\sigma^{(i)}
 =N(\frakG_{i-1}\mid\frakM_{i-1})=N(\frakg_{i-1}\mid\frakg_i).
\end{equation}
Wie (3) zeigt, wird $R^{(i)}$ bei Adjunktion von $x_{i+1},\ldots,x_n$ gleich der Determinante der Übergangssubstitution von $\frakG_{i-1}$ nach $\frakM_{i-1}$, und der Grad von $R^{(i)}$ gibt die Anzahl der in bezug auf $P(u,x_{i+1},\ldots,x_n)$ linear unabhängigen Restklassen von $\frakG_{i-1}$ nach $\frakM_{i-1}$, also von $\frakg_{i-1}$ nach $\frakg_i$ an; $R^{(i)}$ kann als ganz und primitiv in $x_{i+1},\ldots,x_n$ angenommen werden und wird dann regulär in $x_i$. $R^{(i)}$ läßt sich berechnen als größter gemeinsamer Teiler -- im Polynomsinn -- aller $\varrho$-reihigen Determinanten irgendeines stets existierenden Moduls $\frakM^*_{i-1}$ vom Range $\varrho$ von der Eigenschaft, daß das Restklassensystem $\frakG^*_{i-1}/\frakM^*_{i-1}$ isomorph wird zu $\frakG_{i-1}/\frakM_{i-1}$, unter $\frakG^*_{i-1}$ den Grundmodul von $\frakM^*_{i-1}$ verstanden (H.-N. Satz VII und § 5).

\textbf{6. Elementarteilerform und Norm (Resultantenform) der Ideale.} Das Produkt $E_\frakm=E^{(1)}E^{(2)}\cdots E^{(n)}$ aller höchsten Elementarteiler wird als Elementarteilerform, das Produkt $R_\frakm=R^{(1)}R^{(2)}\cdots R^{(n)}$ aller Einzelnormen als Norm (Resultantenform)\footnote{Bei H.-N. wird nur die Bezeichnung Resultantenform gebraucht; die Bezeichnung Norm entspricht den arithmetischen Eigenschaften.} von $\frakm$ bezeichnet. Elementarteilerform und Norm werden durch das Ideal teilbar; die Norm wird durch die Elementarteilerform, eine Potenz der letzteren durch die Norm teilbar. Allgemeiner gilt noch:
\begin{align}
 E^{(i)}\frakg_{i-1}&\equiv0\pmod{\frakg_i},
& E^{(n)}\cdots E^{(i)}\frakg_{i-1}&\equiv0\pmod{\frakm},\notag\\
 R^{(i)}\frakg_{i-1}&\equiv0\pmod{\frakg_i},
& R^{(n)}\cdots R^{(i)}\frakg_{i-1}&\equiv0\pmod{\frakm}, \tag{4}
\end{align}
(H.-N. Satz VIII), und weiter: Ist $\frakm$ teilbar durch $\frakn$, und stimmen die Normen $R_\frakm$ und $R_\frakn$ überein, so stimmen auch die Ideale $\frakm$ und $\frakn$ überein (H.-N. Satz IX). Beachtet man, daß für einen Teiler von $\frakm$ aus dem Übereinstimmen von $R^{(1)},R^{(2)},\ldots$ das Übereinstimmen der Grundideale $\frakg_1,\frakg_2,\ldots$ folgt und daß stets $\frakg_0=\frako$ ist, so kommt ganz entsprechend: Ist $\frakn$ ein echter Teiler von $\frakm$, so wird die erste Einzelnorm von $\frakn$, die von der entsprechenden von $\frakm$ verschieden ist, ein echter Teiler.\footnote{Dagegen können spätere Faktoren von $R_\frakn$ Vielfache der entsprechenden Faktoren von $R_\frakm$ werden; z. B. immer, wenn $\frakm$ Primideal ist und $\frakn$ nicht das Einheitsideal.} Enthält $\frakm$ ein Polynom $G^{(i)}\ne0$, so werden nach Definition $\frakg_0,\frakg_1,\ldots,\frakg_{i-1}$ gleich dem Einheitsideal; und somit werden $R^{(1)},\ldots,R^{(i-1)}$ nach der in 5. angegebenen Eigenschaft der Einzelnormen, die Anzahl linear unabhängiger Restklassen zu bestimmen, gleich der Einheit; folglich werden auch $E^{(1)},\ldots,E^{(i-1)}$ gleich der Einheit.

\textbf{7. Zerlegung der Einzelnormen und Elementarteiler; Komponenten der Grundideale.} Unter einer Primfunktion wird ein in bezug auf $P(u)$ irreduzibles Polynom verstanden, unter Primärfunktion die Potenz einer Primfunktion; eine eigentliche Primärfunktion ist eine solche, die nicht zugleich Primfunktion ist, wo es sich also um höhere als erste Potenz handelt. Eine Zerlegung eines Polynoms in größte primäre Faktoren ist eine solche, wo jeder Faktor Primärfunktion ist, das Produkt irgend zweier Faktoren aber nicht mehr primär.\footnote{Die Definition von Primfunktion und Primärfunktion ließe sich auch in genauer Analogie mit 11. fassen, indem man nur das Wort Ideal durch Polynom ersetzt. Die größten primären Faktoren sind das Analogon zu den größten primären Komponenten der Zerlegung in 12.} Der Zerlegung der Einzelnorm $R^{(i)}=N(\frakg_{i-1}\mid\frakg_i)$ in größte primäre Faktoren $Q_\nu^{(i)}$ entspricht eine Darstellung von $\frakg_i$ als kleinstes gemeinsames Vielfaches
\[
 \frakg_i=[\frakr_1,\frakr_2,\ldots,\frakr_s],
\]
derart, daß $Q_\nu^{(i)}\frakg_{i-1}\equiv0\pmod{\frakr_\nu}$ wird; $\frakr_\nu$ besitzt dabei $\frakg_{i-1}$ als Grundideal $(i-1)$-ter Stufe, und stimmt mit seinem Grundideal $i$-ter Stufe überein; die $\frakr_\nu$ werden als die Komponenten der Grundideale bezeichnet. Der höchste Elementarteiler von $\frakr_\nu$ wird gleich dem größten gemeinsamen Teiler -- im Polynomsinn -- von $Q_\nu^{(i)}$ und $E^{(i)}$, also gleich einem größten primären Faktor von $E^{(i)}$. Die $\frakr_\nu$ sind eindeutig bestimmt durch ihr oben angegebenes Verhalten in bezug auf die Grundideale und durch die Forderung, daß entweder Einzelnorm oder höchster Elementarteiler ein größter primärer Faktor von $R^{(i)}$ bzw. $E^{(i)}$ werden soll (H.-N. Satz X).

\textbf{8. Dimensionsbegriff.} Dem Ideal $\frakm$ wird die Höchstdimension $n-i$ zugeschrieben, wenn $R^{(i)}$ die erste von der Einheit verschiedene Einzelnorm ist -- oder was nach 6. damit gleichbedeutend ist, wenn $\frakg_i$ das erste vom Einheitsideal verschiedene Grundideal ist; dem Einheitsideal $\frako$ wird die Dimension $-1$ zugeschrieben. Ist nur eine von der Einheit verschiedene Einzelnorm $R^{(i)}$ vorhanden, wird also $\frakg_i=\frakg_{i+1}=\cdots=\frakm$, so wird die Höchstdimension auch als Dimension von $\frakm$ bezeichnet. Enthält $\frakm$ ein Polynom $G^{(i)}\ne0$, so ist es höchstens von der Höchstdimension $n-i$, wie die Bemerkung am Schluß von 6. zeigt. Ist $\frakm$ von der Höchstdimension $n-i$, $\frakn$ ein Teiler $\frakm$, so ist $\frakn$ also höchstens von der Höchstdimension $n-i$.

\textbf{9. Norm (Resultantenform) und Nullstellen.} Unter Nullstellen eines Ideals $\frakm$ werden alle einem geeigneten algebraischen Erweiterungskörper von $P(u;x_{i+1},\ldots,x_n)$ angehörigen Wertsysteme von $x_1,\ldots,x_i$ verstanden ($i=1,2,\ldots,n$), für welche -- unter Adjunktion von $x_{i+1},\ldots,x_n$ zum Koeffizientenbereich -- alle Polynome aus $\frakm$ verschwinden. Jedem Faktor $R^{(i)}$ der Norm (Resultantenform) entsprechen bei dieser Adjunktion so viele Nullstellen von $\frakm$, wie der Grad von $R^{(i)}$ angibt; und es läßt $R^{(i)}$, geschrieben in bilinearen Verbindungen $w_{\mu\nu}$ der $u_{\mu\nu}$ und weiterer Unbestimmten $v_{\mu\nu}$, die explizite Zerlegung zu:
\begin{equation}
 R^{(i)}=\prod_\nu (z_i-z_i^{(\nu)})
 =\prod_\nu (v_{i1}x_1+\cdots+v_{in}x_n - z_i^{(\nu)}), \tag{5}
\end{equation}
wo $x_1^{(\nu)},\ldots,x_i^{(\nu)}$ je ein System zusammengehöriger Nullstellen bedeutet (H.-N. Satz XIII; eine neue Begründung wird in § 3 gegeben werden). Unter den Nullstellen eines Ideals von der Höchstdimension $n-i$ gibt es also solche, die von $(n-i)$ Parametern abhängen, aber keine, die von mehr Parametern abhängen; damit ist aber die Übereinstimmung des unter 8. gegebenen Dimensionsbegriffs mit der gewöhnlich üblichen Fassung gezeigt.

\textbf{10. Der Ringbereich der allgemeinen Idealtheorie.} Der zugrunde gelegte Bereich sei ein kommutativer Ring, in dem der Satz von der endlichen Kette gilt: Jede Kette von Idealen $\fraka_1,\fraka_2,\ldots,\fraka_k,\ldots$, wo jedes $\fraka_i$ ein echter Teiler -- Teiler im Idealsinn -- des unmittelbar vorangehenden ist, bricht im Endlichen ab. Diese Forderung ist äquivalent mit der anderen, daß jedes Ideal eine Idealbasis besitzt (Idealtheorie, Satz I), ist also insbesondere für den Polynombereich erfüllt. In den folgenden Nummern 11., 12., 13. wird dieser allgemeine Ringbereich zugrunde gelegt.

\textbf{11. Primideale und Primärideale.} Ein Ideal $\frakp$ heißt Primideal, wenn aus der Teilbarkeit eines Produktes durch $\frakp$ die Teilbarkeit mindestens eines Faktors folgt; $\frakq$ heißt Primärideal -- oder primär --, wenn aus der Teilbarkeit eines Produktes durch $\frakq$ die Teilbarkeit eines Faktors oder einer Potenz jedes Faktors folgt. In Zeichen:
\[
 a\not\equiv0\pmod{\frakp},\quad b\not\equiv0\pmod{\frakp}\quad\hbox{folgt}\quad ab\not\equiv0\pmod{\frakp};
\]
\[
 a\not\equiv0\pmod{\frakq},\quad b^x\not\equiv0\pmod{\frakq}\ \hbox{für jede Potenz }x\quad\hbox{folgt}\quad ab\not\equiv0\pmod{\frakq}.
\]
Zu jedem Primärideal $\frakq$ gibt es ein und nur ein zugehöriges Primideal $\frakp$, das Teiler von $\frakq$ ist und von dem eine Potenz durch $\frakq$ teilbar wird, $\frakp\ne\frako$, $\frakp^\rho\equiv0\pmod{\frakq}$; dabei besteht $\frakp$ aus allen Ringelementen, von denen eine Potenz durch $\frakq$ teilbar wird. Wie die Definition zeigt, ist jedes Primideal zugleich Primärideal; unter eigentlichen Primäridealen werden solche verstanden, die nicht zugleich Primideal sind, für die also $\rho>1$ wird.

\textbf{12. Der Zerlegungssatz.} Eine Darstellung $\fraka=[\frakq_1,\ldots,\frakq_s]$ eines Ideals als kleinstes gemeinsames Vielfaches heißt eine kürzeste, wenn kein $\frakq$ im kleinsten gemeinsamen Vielfachen der übrigen -- in seinem Komplement -- aufgeht; sie heißt eine solche durch größte primäre Komponenten, wenn jedes $\frakq$ primär ist, aber das kleinste gemeinsame Vielfache irgend zweier $\frakq$ nichtprimär ist. Jedes Ideal läßt eine kürzeste Darstellung durch endlich viele größte primäre Komponenten zu; bei zwei verschiedenen solchen Darstellungen stimmen die Anzahl der Komponenten und die zugehörigen Primideale, die alle voneinander verschieden sind, überein. Die derart eindeutig bestimmten Primideale sollen als die Primideale oder zugehörigen Primideale des Ideals bezeichnet werden (Idealtheorie, Satz IX).

\textbf{13. Isolierte Komponenten der Zerlegung.} Ein Ideal $\frakr=[\frakq_1,\ldots,\frakq_i]$ wird isolierte Komponente der Zerlegung von $\fraka$ genannt, wenn die $\frakq_i$ in mindestens einer kürzesten Darstellung von $\fraka$ durch größte primäre Komponenten auftreten und die Bedingung erfüllen, daß keines ihrer zugehörigen Primideale in einem der übrigen Primideale von $\fraka$ aufgeht. Die isolierten Komponenten der Zerlegung sind eindeutig durch ihre Primideale bestimmt; insbesondere sind also die isolierten größten primären Komponenten eindeutig bestimmt (Idealtheorie, Satz XIII).

\textbf{14. Charakteristik, vollkommene und unvollkommene Körper, Erweiterung erster Art.} Ein Körper $\Omega$ heißt von der Charakteristik Null, wenn der in $\Omega$ enthaltene, aus der Einheit abgeleitete Primkörper vom Typus des Körpers der rationalen Zahlen ist; von der Charakteristik $p$, wenn dieser Primkörper vom Typus des Restklassensystems nach einer Primzahl $p$ ist. Ein Körper $\Omega$ heißt vollkommen, wenn jede Primfunktion mit Koeffizienten aus $\Omega$ in einem geeigneten Erweiterungskörper in getrennte Linearfaktoren zerfällt, im entgegengesetzten Falle unvollkommen. Jeder Körper von der Charakteristik Null ist vollkommen; ein Körper von der Charakteristik $p$ aber dann und nur dann, wenn er neben jedem Element auch dessen $p$-te Wurzel enthält. Eine Primfunktion in einem unvollkommenen Körper ist eine ganze Funktion $r$-ten Grades von $x^{p^f}$, und zerfällt in einem geeigneten Erweiterungskörper in $p^f$-te Potenzen von $r$ verschiedenen Linearfaktoren; $f$ wird als Exponent bezeichnet. Eine Primfunktion heißt von erster Art, wenn sie in einem geeigneten Erweiterungskörper in getrennte Linearfaktoren zerfällt, der Exponent $f$ Null wird; eine Erweiterung heißt erster Art, wenn jedes Element erster Art, d. h. Nullstelle einer Primfunktion erster Art wird. Jede Erweiterung, die durch Adjunktion eines Elementes erster Art entsteht, ist erster Art; bei vollkommenen Körpern gibt es nur Erweiterungen erster Art (Steinitz, § 11 und § 13).

\textbf{15. Reduzierter Grad und Exponent einer endlichen Erweiterung.} Ist $\mathcal K$ eine endliche Erweiterung eines unvollkommenen Körpers $\Omega$, so ist in $\mathcal K$ ein Teilkörper $\mathcal K_0$ vom Grade $r$ enthalten, der erster Art in bezug auf $\Omega$ wird und der aus allen und nur den Elementen erster Art aus $\mathcal K$ besteht. Die $p^f$-te Potenz jedes Elementes aus $\mathcal K$ gehört zu $\mathcal K_0$, und es gibt Elemente in $\mathcal K$, die genau vom Grade $rp^f$ sind. $r$ wird als reduzierter Grad, $f$ als Exponent der endlichen Erweiterung bezeichnet (Steinitz, § 14, 1 und § 14, 5).

\subsection*{§ 2. Elementarteilerform und Norm der Primideale und Primärideale. Teiler der Primideale}

Es handelt sich von jetzt an durchweg um den in § 1, 1. definierten Polynombereich. Vorerst ist zu zeigen, daß man sich auch bei Primidealen und Primäridealen auf transformierte Ideale beschränken kann, nach

\textbf{Hilfssatz I.} Das transformierte Ideal $\frakp$ eines Primideals $\bar{\frakp}$ in $\frakS$ wird Primideal, das transformierte $\frakq$ eines Primärideals $\bar{\frakq}$ in $\frakS$ wird Primärideal. M. a. W. die Eigenschaft, Primideal oder Primärideal zu sein, bleibt bei Adjunktion der $u_{\mu\nu}$ zu $P$ erhalten.

Sei nämlich vorausgesetzt: $\fraka\not\equiv0\pmod{\frakp}$, $\frakb\not\equiv0\pmod{\frakp}$, wo $\fraka$ und $\frakb$ nicht notwendig transformierte Ideale zu sein brauchen; sei etwa $a(x)\not\equiv0\pmod{\frakp}$, $b(x)\not\equiv0\pmod{\frakp}$, wo $a(x)$ und $b(x)$ Polynome aus $\fraka$ und $\frakb$ bedeuten. Durch Umkehrung von (1) nach (1') kommt -- unter $T$ ohne Beschränkung der Allgemeinheit Potenzprodukte der $t_{\mu\nu}$ verstanden --
\[
 a(x)=\sum T_i a_i(y),\qquad b(x)=\sum T_i' b_i(y),
\]
und es muß mindestens ein $a_i(y)$ und ein $b_i(y)$ nicht durch $\bar{\frakp}$ teilbar sein. Seien etwa $a_h(y)\not\equiv0\pmod{\bar{\frakp}}$, $b_k(y)\not\equiv0\pmod{\bar{\frakp}}$ die jeweils höchsten Glieder bei irgendeiner lexikographischen Anordnung der $T$; dann wird auch $a_h(y)b_k(y)\not\equiv0\pmod{\bar{\frakp}}$ und somit kommt $a(x)b(x)\not\equiv0\pmod{\frakp}$ und also $\fraka\frakb\not\equiv0\pmod{\frakp}$, womit $\frakp$ als Primideal nachgewiesen ist. Der Beweis beruht offenbar darauf, daß das System der Restklassen nach einem Primideal einen Ring ohne Nullteiler bildet, eine Eigenschaft, die bei Adjunktion von Unbestimmten erhalten bleibt.

Sei entsprechend $\fraka\not\equiv0\pmod{\frakq}$, $\frakb^x\not\equiv0\pmod{\frakq}$ für jede Potenz $x$; dann folgt aus der Existenz der Idealbasis, daß ein $a(x)$ aus $\fraka$ und ein $b(x)$ aus $\frakb$ existieren muß, so daß $a(x)\not\equiv0\pmod{\frakq}$, $b(x)^x\not\equiv0\pmod{\frakq}$ für jede Potenz $x$ wird; denn bei der entgegengesetzten Annahme würde eine Potenz von $\frakb$ durch $\frakq$ teilbar.

Setzt man $a(x)b(x)=c(x)$, so mögen $\fraka^*,\frakb^*,\frakc^*$ die bzw. aus den Koeffizienten $a_i(y),b_j(y),c_k(y)$ von $a(x),b(x),c(x)$ abgeleiteten Ideale bedeuten. Dann wird $\frakb^{*x}\not\equiv0\pmod{\bar{\frakq}}$ für jede Potenz $x$, da sonst $b(x)^x$ durch $\frakq$ teilbar würde; wegen $\fraka^*\not\equiv0\pmod{\bar{\frakq}}$ muß daher auch $\fraka^*\frakb^{*\lambda}\not\equiv0\pmod{\bar{\frakq}}$ für jede Potenz $\lambda$ werden. Nun existiert aber nach dem Dedekind-Mertensschen Satze\footnote{Vgl. H.-N., S. 63 und Anm. 10).} ein Exponent $\lambda$ derart, daß $\fraka^*\frakb^{*\lambda}\equiv \frakc^*\frakb^{*\lambda-1}$ wird; also muß $\frakc^*\not\equiv0\pmod{\bar{\frakq}}$ werden. Das ergibt aber $c(x)\not\equiv0\pmod{\frakq}$ und somit $\fraka\frakb\not\equiv0\pmod{\frakq}$, womit $\frakq$ als primär nachgewiesen ist.

Auch bei der Darstellung eines Ideals als kleinstes gemeinsames Vielfaches kann man sich auf transformierte Ideale beschränken, nach

\textbf{Hilfssatz II.} Aus einer Darstellung $\barfrakm=[\frakc_1,\ldots,\frakc_s]$ in $\frakS$ folgt eine Darstellung $\frakm=[\frakc_1',\ldots,\frakc_s']$ in $\frakS'$, wo $\frakc_i'$ die transformierten von $\frakc_i$ bedeuten. Insbesondere können also in jeder kürzesten Darstellung durch größte primäre Komponenten diese als transformiert angenommen werden.

Es wird nämlich $\frakm$ durch $[\frakc_1',\ldots,\frakc_s']$ teilbar, da jedes Polynom $f(x)$ aus $\frakm$ -- eventuell nach Multiplikation mit einer geeigneten Potenz von $U(u)$ -- von der Form $\sum U_i(u)f_i(y)$ ist, wo jedes $f_i(y)$ als Polynom aus $\barfrakm$ nach Voraussetzung durch alle $\frakc_i$ teilbar ist. Ist umgekehrt $f(x)$ durch jedes $\frakc_i'$ teilbar, so wird es von obiger Form, und somit durch $\frakm$ teilbar. Geht man also von einer Zerlegung in größte primäre Komponenten in $\frakS$ aus, so ergibt dies eine Zerlegung $\frakm=[\frakq_1',\ldots,\frakq_s']$ in $\frakS'$; hier sind die $\frakq_i'$ als transformierte von primären Idealen nach Hilfssatz I primär, und zwar handelt es sich um größte primäre Komponenten, da verschiedenen zugehörigen Primidealen $\bar{\frakp}$ in $\frakS$ auch verschiedene $\frakp$ in $\frakS'$ entsprechen.

Es handelt sich jetzt um eine erste Charakterisierung von Primidealen und Primäridealen durch Elementarteilerform und Norm, noch ohne vollständige Trennung von prim und primär, und zwar gilt in genauer Analogie zu der Idealtheorie in algebraischen Zahlkörpern\footnote{Als höchsten Elementarteiler der Ideale in algebraischen Zahlkörpern hat man die kleinste durch das Ideal teilbare ganze rationale Zahl aufzufassen, also insbesondere die durch ein Primideal teilbare Primzahl oder die durch ein Primärideal (Primidealpotenz) teilbare Primzahlpotenz. In der Tat ist ja jedes Ideal $\fraka^*$ zugleich ein Modul $A^*$ aus Linearformen in den Elementen $\omega_1,\ldots,\omega_n$ einer Körperbasis, $(\omega_1,\ldots,\omega_n)$ Grundmodul von $A^*$; die kleinste durch $\fraka^*$ teilbare ganze rationale Zahl wird also der höchste Elementarteiler dieses Moduls $A^*$; die Norm von $\fraka^*$ aber wird das Produkt aller Elementarteiler von $A^*$.}

\textbf{Satz I.} Die Elementarteilerform eines Primideals wird gleich einer Primfunktion, die Norm gleich einer Potenz dieser Primfunktion. Bei Primäridealen werden Elementarteilerform und Norm gleich Primärfunktionen, und zwar gleich Potenzen derselben Primfunktion, die selbst Elementarteilerform des zugehörigen Primideals ist. Bei Primidealen und Primäridealen fällt die Höchstdimension mit der Dimension schlechthin zusammen; diese Dimension stimmt für Primärideal und zugehöriges Primideal überein.

Satz I ist offenbar erfüllt für das Einheitsideal, dessen Elementarteilerform und Norm gleich der Einheit werden. Sei jetzt das Primideal $\frakp$ von der Höchstdimension $(n-i)$; nach (4), § 1, 6. kommt
\[
 E^{(i)}\frakg_{i-1}\not\equiv0\pmod{\frakp},\quad\hbox{aber}\quad
 E^{(i)}\frakg_i\equiv0\pmod{\frakp},
\]
da $\frakp$ sonst nach § 1, 8. höchstens von der Höchstdimension $n-i-1$ würde. Somit wird $\frakg_i\equiv0\pmod{\frakp}$, also $\frakp$ mit seinem Grundideal $i$-ter Stufe und folglich auch mit $\frakg_{i+1},\frakg_{i+2},\ldots$ identisch. Es werden also $R^{(i+1)},\ldots,R^{(n)}$ gleich der Einheit, mithin auch $E^{(i+1)},E^{(i+2)},\ldots$ als Teiler von $R^{(i+1)},R^{(i+2)},\ldots$; und somit wird die Elementarteilerform gleich $E^{(i)}$, das gleich einer Primfunktion $P^{(i)}$ werden muß. Denn nach § 1, 5. und unter Berücksichtigung, daß $\frakg_{i-1}$ gleich dem Einheitsideal wird, ist $E^{(i)}$ definiert als größter gemeinsamer Teiler -- im Polynomsinn -- aller durch $\frakp$ teilbaren $B^{(i)}$; aus $E^{(i)}=E_1^{(i)}E_2^{(i)}\equiv0\pmod{\frakp}$ würde also folgen, daß $E_1^{(i)}$ in einem seiner Faktoren $E_1^{(i)}$ oder $E_2^{(i)}$ aufgehen muß. Da aber eine Potenz von $E^{(i)}$ durch $R^{(i)}$ teilbar wird, wird $R^{(i)}$ gleich einer Potenz -- die auch die erste sein kann -- dieser Primfunktion $P^{(i)}$; zugleich wird $R^{(i)}$ Norm von $\frakp$. Da nur eine von der Einheit verschiedene Einzelnorm $R^{(i)}$ auftritt, stimmt schließlich die Höchstdimension von $\frakp$ mit der Dimension schlechthin überein.

Sei entsprechend das Primärideal $\frakq$ von der Höchstdimension $n-i$; dann kommt wie oben
\[
 (E^{(i)}\frakg_{i-1})^x\not\equiv0\pmod{\frakq},\quad\hbox{aber}\quad
 (E^{(i)}\frakg_i)^x\equiv0\pmod{\frakq}
\]
für jede Potenz $x$; und folglich wird wie oben $\frakq=\frakg_i$, seine Elementarteilerform gleich $E^{(i)}$, seine Norm gleich $R^{(i)}$, die Höchstdimension gleich der Dimension schlechthin. Diese Dimension stimmt mit der des zugehörigen Primideals überein; denn aus $\frakp^\rho\equiv0\pmod{\frakq}$, $\frakq\equiv0\pmod{\frakp}$ folgt, daß die Dimension $(n-i)$ von $\frakp$ höchstens gleich der von $\frakq$, die von $\frakq$ höchstens gleich der von $\frakp^\rho$ ist; aus $(P^{(i)})^\rho\equiv0\pmod{\frakp^\rho}$, wo $P^{(i)}$ die Elementarteilerform von $\frakp$ bedeutet, ergibt sich letztere Höchstdimension als höchstens $n-i$, und somit wird $n-i$ die Dimension von $\frakp$ und $\frakq$. Aus $(P^{(i)})^\rho\equiv0\pmod{\frakq}$ folgt ferner, daß $E^{(i)}$ -- als größter gemeinsamer Teiler aller $B^{(i)}$ aus $\frakq$ -- eine Potenz von $P^{(i)}$ wird, die auch die erste sein kann, und daß somit das gleiche für $R^{(i)}$ gilt, womit Satz I in allen Teilen bewiesen ist.

Eine Charakterisierung der Primideale vermöge des Dimensionsbegriffs wird gegeben durch

\textbf{Satz II.} Ein Primideal besitzt keinen echten Teiler gleicher Höchstdimension; und umgekehrt wird jedes solche Ideal Primideal.

Der Beweis beruht auf

\textbf{Hilfssatz III.} Besitzt ein Primideal $\frakp$ aus Polynomen mit Koeffizienten aus einem Körper $\Omega$ nur endlichviele in bezug auf $\Omega$ linear unabhängige Restklassen, so hat $\frakp$ keinen vom Einheitsideal verschiedenen echten Teiler; m. a. W. das Restklassensystem nach $\frakp$ bildet einen Körper.

Die Voraussetzung besagt, daß es eine endliche Zahl $k$ gibt derart, daß zwischen je $(k+1)$ Restklassen eine lineare Abhängigkeit mit Koeffizienten aus $\Omega$ besteht -- genauer mit Koeffizienten aus dem durch alle Elemente von $\Omega$ repräsentierten Restklassenkörper $(\Omega)$, daß aber mindestens ein System von $k$ in bezug auf $\Omega$ linear unabhängigen Restklassen existiert; daraus folgt sofort, daß sich alle Restklassen linear durch ein beliebiges System von $k$ linear unabhängigen ausdrücken lassen. Sei jetzt $\fraka$ ein echter Teiler von $\frakp$, $a(z)\not\equiv0\pmod{\frakp}$ ein Polynom aus $\fraka$; aus $c_1f_1(z)+\cdots+c_kf_k(z)\equiv0\pmod{\frakp}$ folgt dann auch
\[
 c_1a(z)f_1(z)+\cdots+c_ka(z)f_k(z)\equiv0\pmod{\frakp};
\]
das heißt aber, daß neben den Restklassen von $f_1,\ldots,f_k$ auch die von $af_1,\ldots,af_k$ ein System von $k$ linear unabhängigen bilden, durch die sich also insbesondere die Einheitsklasse linear darstellen läßt. Es wird also $\fraka$ gleich dem Einheitsideal; anders ausgedrückt, das Restklassensystem bildet einen Körper, da es zu $a(z)\not\equiv0\pmod{\frakp}$ immer ein $f(z)$ gibt, so daß $1\equiv a(z)f(z)\pmod{\frakp}$ wird.

Zum Beweise des ersten Teiles von Satz II ist jetzt nur zu bemerken, daß jedes Primideal vor der Dimension $(n-i)$ -- allgemeiner jedes Ideal von der Höchstdimension $(n-i)$ -- nur endlich viele in bezug auf $P(u;x_{i+1},\ldots,x_n)$ linear unabhängige Restklassen besitzt, nämlich (§ 1, 5.) so viele wie der Grad von $R^{(i)}$ angibt, da $\frakg_{i-1}$ hier gleich dem Einheitsideal wird. Jeder echte Teiler $\fraka$ von $\frakp$ geht also bei Adjunktion von $x_{i+1},\ldots,x_n$ zu $P(u)$ in das Einheitsideal über; anders ausgedrückt, das Grundideal $i$-ter Stufe von $\fraka$ wird gleich dem Einheitsideal; das heißt, daß die Höchstdimension von $\fraka$ höchstens gleich $n-i-1$ wird. Damit die Aussage auch für ein nichttransformiertes Ideal $\fraka$ als Teiler gilt, hat man nur nach § 1, 2. zu einem neuen transformierten Bereich überzugehen.

Besitze jetzt umgekehrt $\frakp$ keinen echten Teiler gleicher Höchstdimension $n-i$; und sei $\fraka\not\equiv0\pmod{\frakp}$, $\frakb\not\equiv0\pmod{\frakp}$, wo $\fraka$ und $\frakb$ wieder als transformierte Ideale angenommen sind. Dann kommt nach Voraussetzung, da $(\fraka,\frakp)$ und $(\frakb,\frakp)$ als größte gemeinsame Teiler -- im Idealsinn -- echte Teiler von $\frakp$ werden:
\[
 E^{(n)}\cdots E^{(i+1)}\frakg_i\equiv0\pmod{(\fraka,\frakp)},
\qquad
 E^{(n)}\cdots E^{(i+1)}\frakg_i\equiv0\pmod{(\frakb,\frakp)}.
\]
Das ergibt aber
\[
 E^{(n)}\cdots E^{(i+1)}\frakg_i\equiv0\pmod{(\fraka\frakb,\frakp)},
\]
und somit folgt notwendig $\fraka\frakb\not\equiv0\pmod{\frakp}$, da sonst $\frakp$ von geringerer Höchstdimension als $n-i$ wäre. Da $\fraka,\frakb$ transformierte Ideale sind, ergibt sich also $\frakp$ und nach Hilfssatz I auch $\bar{\frakp}$ als Primideal, womit Satz II bewiesen ist.

\subsection*{§ 3. Nullstellen der Primideale und Primärideale}

Den Übergang zu einer direkten Begründung der Theorie der Nullstellen (§ 1, 9.), vorerst für Primideale, bildet der eine Erweiterung von Hilfssatz III darstellende und für Primideale in beliebigen Ringen gültige

\textbf{Hilfssatz IV.} Das Restklassensystem nach jedem vom Einheitsideal verschiedenen Primideal läßt sich durch Quotientenbildung zu einem Körper erweitern, dem Restklassenkörper des Primideals.

Das Restklassensystem nach einem beliebigen Ideal bildet einen Ring, da Summe, Differenz und Produkt wieder dem System angehören, und assoziatives, kommutatives und distributives Gesetz beim Übergang zu den Restklassen erhalten bleiben. Handelt es sich speziell um Primideale, so besitzt dieser Ring keine Nullteiler; denn die Definition der Primideale (§ 1, 11.) sagt aus, daß das Produkt nicht verschwindender Restklassen ebenfalls nicht verschwindet. Ist das Primideal vom Einheitsideal verschieden, so enthält dieser Ring mindestens ein vom Nullelement verschiedenes Element, und läßt sich daher durch Quotientenbildung -- Adjunktion von Elementenpaaren -- zum Körper erweitern;\footnote{Steinitz, § 3. Die Annahme eines von Null verschiedenen Elementes im ursprünglichen Integritätsbereich genügt, da dann die Existenz der Einheit durch Quotientenbildung gesichert ist, und somit der ursprüngliche Bereich ein Teilbereich des Körpers wird; Steinitz setzt Existenz der Einheit im Integritätsbereich voraus.} somit ist der Restklassenkörper definiert.

Dem folgenden vorausgeschickt sei die durchweg festgehaltene

\textbf{Bezeichnung:} Restklassen sollen allgemein durch Klammersetzen irgendeines Elementes der Klasse bezeichnet werden, $(x),\ldots,(a(x)),\ldots$; ebenso sollen Körper aus Restklassen durch Klammern bezeichnet werden. Insbesondere bezeichne $(P)$, $(P(u))$, $(P(u,x_{i+1},\ldots,x_n))$ die Restklassenkörper, die aus allen und nur den Restklassen bestehen, die durch Elemente aus $P$, $P(u)$, $P(u,x_{i+1},\ldots,x_n)$ repräsentiert werden können.

Weiter vorausgeschickt sei noch das folgende: Ein Körper $\mathcal R$ heißt vom algebraischen Rang (Transzendenzgrad) $k$ in bezug auf einen Grundkörper $\Omega$, wenn $\mathcal R$ mindestens ein System von $k$ in bezug auf $\Omega$ algebraisch unabhängigen Größen -- transzendenten Größen -- enthält, wenn aber je $k+1$ Größen aus $\mathcal R$ algebraisch abhängig in bezug auf $\Omega$ sind. Läßt sich $\mathcal R$ als algebraischer Erweiterungskörper eines Teilkörpers $\Omega(z_1,\ldots,z_s)$ darstellen, wo die $z$ algebraisch unabhängig -- transzendent -- in bezug auf $\Omega$ sind, so wird notwendig $s$ gleich $k$;\footnote{Für einen bei beliebiger Charakteristik gültigen Beweis dieser bekannten Tatsache vgl. Steinitz, § 22, 9. -- Bei Adjunktion der $u$ kann der algebraische Rang nicht zunehmen, da es sich um rationale Verbindungen der ursprünglichen Elemente mit Koeffizienten aus $\Omega(u)$ handelt; nicht abnehmen, da jede Relation zwischen Elementen aus $\mathcal R$ mit Koeffizienten aus $\Omega(u)$ in endlich viele Relationen mit Koeffizienten aus $\Omega$ zerfällt.} ebenso wird der algebraische Rang von $\mathcal R(u)$ in bezug auf $\Omega(u)$ gleich $k$, wenn $u$ nicht in $\mathcal R$ enthaltene Unbestimmte bedeuten.

Die Konstruktion des Nullstellenkörpers geschieht jetzt in voller Analogie mit dem bei Primfunktionen einer Unbestimmten üblichen Weg\footnote{Steinitz, § 6.} nach

\textbf{Satz III.} Dem Restklassenkörper $(\mathcal R)$ eines Primideals $\frakp$ von der Dimension $n-i$ läßt sich isomorph ein Nullstellenkörper $R$ von $\frakp$ zuordnen, der vom algebraischen Rang $n-i$ wird -- von $n-i$ Parametern abhängt -- und insbesondere als endlicher Erweiterungskörper von $P(u,x_{i+1},\ldots,x_n)$ aufgefaßt werden kann. Der Isomorphismus zwischen $(\mathcal R)$ und $R$ umfaßt denjenigen zwischen $(P)$ und $P$, $(P(u))$ und $P(u)$ und bei Zugrundelegung der $x_{i+1},\ldots,x_n$ als Parameter auch den zwischen $(P(u,x_{i+1},\ldots,x_n))$ und $P(u,x_{i+1},\ldots,x_n)$. Umgekehrt ist jeder Nullstellenkörper vom algebraischen Rang $n-i$ -- von $n-i$ Parametern abhängend -- dem Restklassenkörper $(\mathcal R)$ isomorph.

Satz III ergibt als Korollar den bekannten Satz: Verschwindet ein Ideal $\fraka$ in allen Nullstellen eines Primideals $\frakp$, so wird $\fraka$ durch $\frakp$ teilbar.

Zum Beweise ist vorerst zu bemerken, daß der Restklassenkörper $(\mathcal R)$ vom algebraischen Rang $n-i$ in bezug auf $(P(u))$ wird. Denn die Klassen $(x_{i+1}),\ldots,(x_n)$ sind in bezug auf $(P(u))$ algebraisch unabhängig, da $\frakp$ als von der Dimension $n-i$ kein von $x_1,\ldots,x_i$ freies Polynom enthält; jede weitere Klasse aber ist von diesen abhängig. Denn aus der Zugehörigkeit der Elementarteilerform $P^{(i)}$ zu $\frakp$ folgt nach § 1, 2. auch für $\lambda=1,2,\ldots,i$ die Zugehörigkeit von Polynomen, die jeweils von $x_\lambda,x_{i+1},\ldots,x_n$ abhängig sind. $(\mathcal R)$ wird also ein endlicher Erweiterungskörper von $(P(u,x_{i+1},\ldots,x_n))$; der Grad in bezug auf diesen Teilkörper wird nach § 1, 5. -- unter Berücksichtigung, daß $\frakg_{i-1}$ gleich dem Einheitsideal, $\frakg_i$ nach Satz I gleich $\frakp$ wird -- gleich dem Grad der Norm.

Um zu einem isomorphen Nullstellenkörper überzugehen, ist weiter zu bemerken, daß $P(u,x_{i+1},\ldots,x_n)$ und $(P(u,x_{i+1},\ldots,x_n))$ isomorph aufeinander bezogen sind dadurch, daß jedem Element aus $P(u,x_{i+1},\ldots,x_n)$ die durch das Element repräsentierte Restklasse zugeordnet wird. Denn vorerst entsprechen sich bei dieser Zuordnung die aus den Polynomen in $u,x_{i+1},\ldots,x_n$ bzw. in $(u),(x_{i+1}),\ldots,(x_n)$ bestehenden Integritätsbereiche isomorph, da in $\frakp$ kein von $x_1,\ldots,x_i$ freies Polynom enthalten ist, verschiedenen Polynomen aus $P(u,x_{i+1},\ldots,x_n)$ also auch verschiedene Klassen entsprechen; aus isomorphen Integritätsbereichen gehen aber durch Quotientenbildung isomorphe Körper hervor. Die isomorphe Beziehung zwischen $P(u,x_{i+1},\ldots,x_n)$ und $(P(u,x_{i+1},\ldots,x_n))$ umfaßt diejenige zwischen $P$ und $(P)$, $P(u)$ und $(P(u))$. Um diesen Isomorphismus zu einem solchen zwischen $(\mathcal R)$ und einem Nullstellenkörper $R$ zu erweitern, seien den Klassen $(x_1),\ldots,(x_i)$ gewisse neue Elemente $\xi_1,\ldots,\xi_i$ isomorph zugeordnet; d. h. jedem Polynom $(a(x))$ entspreche ein Polynom $a(\xi_1,\ldots,\xi_i,x_{i+1},\ldots,x_n)$; Summe und Produkt entspreche Summe und Produkt. Durch Quotientenbildung -- wobei man sich auf Nenner aus dem Teilkörper $(P(u,x_{i+1},\ldots,x_n))$ bzw. $P(u,x_{i+1},\ldots,x_n)$ beschränken kann -- entspricht dann dem Restklassenkörper $(\mathcal R)$ ein Körper $R$, der endliche Erweiterung von $P(u,x_{i+1},\ldots,x_n)$ vom Grade der Norm wird, und der als Nullstellenkörper bezeichnet werden kann; denn die Zuordnung sagt aus, daß $f(\xi_1,\ldots,\xi_i,x_{i+1},\ldots,x_n)$ verschwindet, sobald $f(x)\equiv0\pmod{\frakp}$ wird. An Stelle von $(P(u,x_{i+1},\ldots,x_n))$ könnte bei der ganzen Überlegung auch irgendein anderer, durch Adjunktion von $n-i$ algebraisch unabhängigen Größen entstehender Teilkörper von $(\mathcal R)$ treten.

Umgekehrt ist leicht zu zeigen, daß jeder Nullstellenkörper $R$ vom algebraischen Rang $n-i$ dem Restklassenkörper isomorph wird. Da sich $R$ rational, mit Koeffizienten aus $P(u)$, aus den Größen $\xi_1,\ldots,\xi_i$ ableiten läßt, müssen unter diesen Größen $n-i$ algebraisch unabhängige, also in bezug auf $P(u)$ transzendente Elemente enthalten sein. Insbesondere muß dies für $\xi_{i+1},\ldots,\xi_n$ erfüllt sein, da wegen $P^{(i)}\equiv0\pmod{\frakp}$ wie oben folgt, daß $R$ ein algebraischer Erweiterungskörper von $P(u,\xi_{i+1},\ldots,\xi_n)$ wird. Da nach Voraussetzung $f(\xi)$ verschwindet für jedes Polynom $f(x)$ aus $\frakp$, entspricht jeder Klasse des in $(\mathcal R)$ enthaltenen Polynombereichs der $(x)$ ein und nur ein Element aus $R$, das ein Polynom in den $\xi$ wird; und Summe und Produkt entsprechen Summe und Produkt. Es entsprechen aber verschiedenen Klassen auch verschiedene Elemente aus $R$; denn sonst würde einer von Null verschiedenen Klasse, also nach geeigneter Multiplikation auch einer durch ein Polynom repräsentierten Klasse aus $(P(u,x_{i+1},\ldots,x_n))$ das Element Null in $R$ entsprechen, zwischen den Größen $\xi_{i+1},\ldots,\xi_n$ bestünde entgegen der Voraussetzung eine algebraische Abhängigkeit. Durch diese Zuordnung werden alle in $R$ enthaltenen Polynome der $\xi$ erschöpft; durch Quotientenbildung -- wobei man sich wieder auf Nenner aus $(P(u,x_{i+1},\ldots,x_n))$ bzw. $P(u,x_{i+1},\ldots,x_n)$ beschränken kann -- gehen aber aus diesen isomorphen Integritätsbereichen isomorphe Körper hervor.

Das Ideal $\fraka$ verschwinde jetzt in allen Nullstellen von $\frakp$, also insbesondere auch in allen von $n-i$ Parametern abhängigen. Ist dann $a(x)$ ein beliebiges Polynom aus $\fraka$, so verschwindet also $a(\xi)$; wegen des Isomorphismus zwischen $R$ und $(\mathcal R)$ wird somit $(a(x))$ die Nullklasse, $a(x)$ und somit $\fraka$ durch $\frakp$ teilbar. Damit ist Satz III und Korollar bewiesen.

Um von Satz III zu der in § 1, 9. gegebenen Zerlegung zunächst für Primideale überzugehen, und um genauere Aussagen über die Exponenten machen zu können, bedarf es einer weiteren Untersuchung der Elementarteilerform nach

\textbf{Hilfssatz V.} Ist bei Charakteristik $p$ die Elementarteilerform $E^{(i)}$ eines Primideals oder Primärideals -- allgemeiner eines Ideals $\fraka$, bei dem die Höchstdimension mit der Dimension schlechthin zusammenfällt -- vom Exponenten $f$ in bezug auf $x_i$, also ganze Funktion von $x_i^{p^f}$, so ist sie auch vom Exponenten $f$ in bezug auf $x_{i+1},\ldots,x_n$ und auf die in $E^{(i)}$ auftretenden Unbestimmten $u_{\mu\nu}$ bzw. $t_{\mu\nu}$. Insbesondere ist also bei vollkommenem Körper als Koeffizientenbereich die Elementarteilerform $P^{(i)}$ eines Primideals stets vom Exponenten Null in bezug auf $x_i$, also eine Primfunktion erster Art.

Dazu ist zu bemerken, daß $P(u,x_{i+1},\ldots,x_n)$ unvollkommen wird, sobald $P$ ein vollkommener Körper von der Charakteristik $p$ ist, so daß aus § 1, 14. nicht direkt folgt, daß $P^{(i)}$ erster Art wird. Sei jetzt $E^{(i)}$ vom Exponenten $f$ in bezug auf $x_i$, enthalte aber eine andere Unbestimmte $x_{i+1}$ zum Exponenten $f'$. Dann gibt es -- durch Vertauschen der $u_{\mu\nu}$ bzw. $t_{\mu\nu}$ -- nach § 1, 2. auch ein von Null verschiedenes durch das Ideal teilbares Polynom $G^{(i)}$, das ganze Funktion von $x_{i+1}^{p^{f'}}$ wird, und das nach Definition der Elementarteilerform durch $E^{(i)}$ teilbar wird. Da aber $G^{(i)}$ von gleichem Grad wie $E^{(i)}$ wird, muß es bis auf eine Größe aus $P(u)$ als Faktor mit $E^{(i)}$ übereinstimmen, und somit wird notwendig $f'$ gleich $f$. Damit wird aber $E^{(i)}$ -- das als ganz und primitiv in den $t_{\mu\nu}$ vorausgesetzt werden darf -- von der Form
\[
 E^{(i)}=\sum c_\rho(t)\,x_i^{\alpha_\rho p^f}x_{i+1}^{\beta_\rho p^f}\cdots x_n^{\gamma_\rho p^f}
       =\sum c_\rho(t)\,X_\rho^{p^f},
\]
wo die $X_\rho$ Potenzprodukte der $x$ bedeuten, und die $X_\rho$ durch $\fraka$ teilbar werden. Daher erhält man ein ebenfalls zu $\fraka$ gehörendes Polynom, wenn in den $X_\rho$ jeder Exponent der einzelnen $t$ durch das größte darin enthaltene Vielfache von $p^f$ ersetzt wird; das kommt aber, wie die obige Darstellung ablesen läßt, darauf hinaus, in $c_\rho(t)$ jeden Exponenten der $t$ durch das größte darin enthaltene Vielfache von $p^f$ zu ersetzen. Durch diese Substitution geht $E^{(i)}$ über in ein Polynom in den $x_i^{p^f},\ldots,x_n^{p^f}$, das nach Definition durch $E^{(i)}$ teilbar wird, und zwar wegen der vorausgesetzten Primitivität auch in bezug auf die $t$; es muß also mit $E^{(i)}$ identisch werden, da der Grad in keinem Argumente zugenommen hat. Ist insbesondere $P$ ein vollkommener Körper, so wird also $E^{(i)}$ als ganze Funktion von $x_i^{p^f}$ eine $p^f$-te Potenz; handelt es sich daher um ein Primideal, so muß notwendig $f$ gleich Null werden; die Elementarteilerform $P^{(i)}$ wird eine Primfunktion erster Art.

Die Zerlegung wird jetzt gegeben durch

\textbf{Satz IV.} Die Elementarteilerform $P^{(i)}$ eines Primideals läßt in dem aus dem Nullstellenkörper abgeleiteten Galoisschen Körper die explizite Zerlegung zu:
\[
 P^{(i)}=\prod_\nu (x_i-t_{i+1}x_{i+1}^{(\nu)}-\cdots-t_nx_n^{(\nu)})^e
 =\prod_\nu (t_i(y_i-y_i^{(\nu)})+\cdots+t_n(y_n-y_n^{(\nu)}))^e,
\]
wo $y_1^{(\nu)},\ldots,y_n^{(\nu)}$ je ein zusammengehöriges von $t_i,\ldots,t_n$ unabhängiges Nullstellensystem der ursprünglichen Unbestimmten $y$ darstellt, verschiedenen $\nu$ verschiedene Linearfaktoren entsprechen und wo der Exponent $e$ bei vollkommenem Körper $P$ den Wert Eins, bei unvollkommenem den Wert $p^f$ ($f>0$) hat. Dadurch ist bei Adjunktion neuer Unbestimmten $v$ die Zerlegung der in bilinearen Verbindungen der $u$ und $v$ an Stelle der $u_{\mu\nu}$ geschriebenen Elementarteilerform mitbestimmt:
\[
 P^{(i)}(w)=\prod_\nu (z_i-z_i^{(\nu)})^e
 =\prod_\nu (v_{i1}(y_1-y_1^{(\nu)})+\cdots+v_{in}(y_n-y_n^{(\nu)}))^e,
\]
wo $e$ die gleiche Bedeutung wie oben hat. Ebenso ist die Zerlegung der Norm eines Primideals oder eines Primärideals mit $\frakp$ als zugehörigem Primideal -- als Potenz von $P^{(i)}$ -- mitgegeben.

Als Korollar aus Satz IV kommt noch: Stimmen zwei Primideale in ihrer Elementarteilerform $P^{(i)}$ überein, so sind sie identisch.

Der Beweis von Satz IV wird auf den Nachweis hinauskommen, daß die Elementarteilerform $P^{(i)}$ mit der Fundamentalgleichung des Erweiterungskörpers $(P(u,x_{i+1},\ldots,x_n),(x_i))$ von $(P(u,x_{i+1},\ldots,x_n))$ identisch ist. Vorerst gehe man -- um Einblick in die Abhängigkeit von den Unbestimmten $u_{\mu\nu}$ bzw. $t_{\mu\nu}$ zu gewinnen -- auf den Restklassenkörper $(\mathcal S)$ von $\frakp$ über, der nach dem bei der Definition des algebraischen Rangs Gesagten vom algebraischen Rang $n-i$ in bezug auf $(P)$ werden muß, wo $(P)$ den durch Elemente aus $P$ repräsentierten Teilkörper bedeutet. Denn $(\mathcal S)$ entsteht durch Adjunktion der Unbestimmten $u$ bzw. $t$ zu einem $(\mathcal R)$ isomorphen Teilkörper, der durch Quotientenbildung aus allen und nur den Klassen hervorgeht, die durch Polynome aus $\frakS$ repräsentiert werden können, wobei $(P)$ und $(P)$ einander entsprechen. Der algebraische Rang von $(\mathcal S)$ in bezug auf $(P(u))$ ist also gleich dem dieses Teilkörpers in bezug auf $(P)$, und somit wegen des Isomorphismus gleich dem von $(\mathcal R)$ in bezug auf $(P)$. Da $(\mathcal S)$ sich rational, mit Koeffizienten aus $(P)$, aus den Klassen $(y_1),\ldots,(y_n)$ ableiten läßt, müssen also unter diesen Klassen $n-i$ algebraisch unabhängige sein, und $(\mathcal S)$ entsteht als endlicher algebraischer Erweiterungskörper dieses durch Adjunktion der $n-i$ Klassen zu $(P)$ erzeugten Teilkörpers.\footnote{Diese Bemerkung zeigt, daß Satz III mit Korollar direkt für den Nullstellenkörper $R$ von $\frakp$ hätte ausgesprochen und bewiesen werden können. Durch den Übergang zum transformierten Ideal wird aber erst erreicht, daß in den größten primären Komponenten gleicher Dimension eines beliebigen Ideals bei den Nullstellen die gleichen Parameter auftreten; und damit ergibt sich erst der Zusammenhang mit Norm und Elementarteilerform eines beliebigen Ideals.}

Adjungiert man nur die in $x_{i+1},\ldots,x_n$ auftretenden Unbestimmten $t_{\mu\nu}$, so entsteht wieder ein Restklassenkörper, der die Klassen $(y_1),\ldots,(y_n)$ und $(x_{i+1}),\ldots,(x_n)$ enthält und endlicher algebraischer Erweiterungskörper von $(P(t_{\mu\nu},x_{i+1},\ldots,x_n))$ wird -- es handelt sich dabei um Körper, die zu Teilkörpern von $(\mathcal S)$ isomorph sind, nicht damit identisch; da bei verschiedenen Unbestimmten die Restklassen zwar durch dieselben Elemente repräsentiert werden, aber nicht identisch sind, da sie nicht die gleichen Elemente enthalten. Bei der Abhängigkeit der Klassen $(y)$ von $(P(t_{\mu\nu},x_{i+1},\ldots,x_n))$ gehen also nur die Unbestimmten $t_{\mu\nu}$ ein; man bilde jetzt, durch Adjunktion von $t_{i1},\ldots,t_{in}$, die Klasse $(x_i)=(t_{i1}y_1+\cdots+t_{in}y_n)$; und es seien $(x_{i\nu})=(t_{i1}y_{1\nu}+\cdots+t_{in}y_{n\nu})$ die $r$ verschiedenen konjugierten Werte, welche in dem in bezug auf $(P(t_{\mu\nu},x_{i+1},\ldots,x_n))$ Galoisschen Körper gelegen sind. Dann wird -- je nachdem es sich um Erweiterungen erster Art handelt oder nicht -- das Produkt über die Faktoren $x_i-(x_{i\nu})$ oder über ihre $p^f$-ten Potenzen eine Primfunktion $(G)$ in bezug auf $(P(t_{\mu\nu},x_{i+1},\ldots,x_n))$; denn da es Elemente vom Grade $r\cdot p^f$ gibt (§ 1, 15.), muß notwendig $(x_i)$ ein solches sein, da jedes Element durch Spezialisierung der $t_{\mu\nu}$ entsteht; $(x_i)$ wird aber Nullstelle von $(G)$. Geht man zu dem isomorphen Nullstellenkörper von $\frakp$ und dem daraus abgeleiteten Galoisschen Körper über, so läßt also $G$ eine Zerlegung in Linearfaktoren $x_i-t_{i1}y_{1\nu}-\cdots-t_{in}y_{n\nu}$ bzw. ihre $p^f$-ten Potenzen zu. Nun wird aber, da $(G)$ für $x_i=(x_i)$ verschwindet, $G(x_i)$ durch $\frakp$ und somit durch $P^{(i)}$ teilbar, und als Primfunktion in bezug auf $P(t_{\mu\nu},x_{i+1},\ldots,x_n)$ und als primitiv in $x_i$ muß $G$ mit $P^{(i)}$ identisch werden; somit ist die erste Zerlegung von Satz IV bewiesen.

Um zu der zweiten Zerlegung überzugehen, ist zu bemerken, daß die Eigenschaft, Erweiterung erster Art zu sein oder nicht, bei Adjunktion von Unbestimmten erhalten bleibt. Bildet man also, nach Adjunktion von $v_1,\ldots,v_n$ und aller $t_{\mu\nu}$, die Klasse $(z)=(v_1x_1+\cdots+v_ix_i)$ und die konjugierten Klassen $(z_\nu)=(v_1x_{1\nu}+\cdots+v_ix_{i\nu})$, so wird das Produkt über alle $z-(z_\nu)$ bzw. ihre $p^f$-ten Potenzen wieder eine Primfunktion $(H)$, die für $z=(z)$ verschwindet. Somit wird $H$ wie oben als Produkt von Linearfaktoren darstellbar und mit der Elementarteilerform des aus $\frakp$ durch Zusammensetzen der Transformation (1) mit $z=v_1x_1+\cdots+v_nx_n$ abgeleiteten Primideals identisch; diese Elementarteilerform entsteht aber -- wegen der gegenseitigen Spezialisierung -- aus $P^{(i)}$ dadurch, daß die $u_{\mu\nu}$ durch gewisse bilineare Verbindungen der $u$ und $v$ ersetzt werden.\footnote{H.-N., § 7.} Mit diesen Zerlegungen ist aber auch die Zerlegung der Norm als Potenz von $P^{(i)}$ mitgegeben,\footnote{Es handelt sich bei Erweiterungen erster Art, also insbesondere bei vollkommenen Körpern, um die erste Potenz; bei unvollkommenen Körpern aber um die $p^f$-te Potenz, wie in § 6, Satz XIII und XIV gezeigt werden wird.} ebenso die von Norm und Elementarteilerform eines Primärideals mit $\frakp$ als zugehörigem Primideal.

Stimmen schließlich zwei Primideale in der Elementarteilerform $P^{(i)}$ überein, so stimmen sie als Folge der obigen Zerlegung in ihren Nullstellen überein, sind also gegenseitig durch einander teilbar und somit identisch.

\end{document}
